\section*{अध्याय \ref{ch:b1:organomagnesium} --- कार्बमैग्नीशियम अभिकर्मक}

\begin{solution}{exo:b1:organomagnesium:1}
\ce{CH3CH2Br + Mg -> CH3CH2MgBr}; \ce{C6H5Br + Mg -> C6H5MgBr} (शुष्क
ईथर में)। C--Br: Br (2.96), C (2.55) से अधिक विद्युत-ऋणात्मक है, कार्बन
धनात्मक। C--Mg: Mg (1.31) कम विद्युत-ऋणात्मक है, कार्बन ऋणात्मक।
\end{solution}

\begin{solution}{exo:b1:organomagnesium:2}
जल: \ce{CH3MgI + H2O -> CH4 + Mg(OH)I}।
एथेनॉल: \ce{CH3MgI + CH3CH2OH -> CH4 + CH3CH2OMgI}।
एथेनोइक अम्ल: \ce{CH3MgI + CH3COOH -> CH4 + CH3COOMgI}।
ड्यूटेरियम ऑक्साइड: \ce{CH3MgI + D2O -> CH3D + Mg(OD)I}।
\end{solution}

\begin{solution}{exo:b1:organomagnesium:3}
मेथेनल से प्रोपेन-1-ऑल (प्राथमिक) मिलता है। एथेनल से ब्यूटेन-2-ऑल (द्वितीयक)।
प्रोपेनोन से 2-मेथिलब्यूटेन-2-ऑल (तृतीयक)। बेंज़ैल्डिहाइड से
1-फ़ेनिलप्रोपेन-1-ऑल (द्वितीयक)। कार्बन डाइऑक्साइड से प्रोपेनोइक अम्ल।
\end{solution}

\begin{solution}{exo:b1:organomagnesium:4}
उसका OH समूह इतना अम्लीय है कि अभिकर्मक बनते ही उसे नष्ट कर दे:
अणु स्वयं को (या अपने पड़ोसियों को) प्रोटॉनित कर देगा, जिससे
ब्यूटेन-1-ऑल का मैग्नीशियम ऐल्कॉक्साइड, \ce{CH3(CH2)3OMgBr}, बनेगा, कोई उपयोगी
अभिकर्मक नहीं। OH को पहले रक्षित करना होगा
(\cref{ch:b1:acetals-protection} किसी ऐल्डिहाइड के लिए यह विचार दिखाता है)।
\end{solution}

\begin{solution}{exo:b1:organomagnesium:5}
C--Mg युग्म एथेनल के कार्बोनिल कार्बन पर आक्रमण करता है, और C=O का $\pi$ युग्म
O पर चला जाता है: \ce{CH3CH(CH3)O^-}, \ce{MgBr+} के साथ; फिर \ce{H3O+}
ऑक्सीजन को एक प्रोटॉन देता है: प्रोपेन-2-ऑल। प्रोपेन-2-ऑल के कार्बिनॉल कार्बन पर
दो मेथिल समूह हैं: वह \omterm{def:b1:stereochemistry-in-depth:stereocentre}{अकिरैल} है। (किसी अन्य R के साथ, जैसे
एथेनल पर एथिलमैग्नीशियम ब्रोमाइड, उत्पाद ब्यूटेन-2-ऑल \omterm{def:b1:stereochemistry-in-depth:stereocentre}{किरैल} है पर
रेसेमिक: समतलीय कार्बोनिल पर दोनों फलकों से समान रूप से आक्रमण होता है।)
\end{solution}

\begin{solution}{exo:b1:organomagnesium:6}
कार्बिनॉल कार्बन पर एक मेथिल और दो एथिल समूह हैं:
ब्यूटेन-2-ओन के साथ एथिलमैग्नीशियम ब्रोमाइड, या पेंटेन-3-ओन के साथ
मेथिलमैग्नीशियम ब्रोमाइड।
\end{solution}

\begin{solution}{exo:b1:organomagnesium:7}
2-ब्रोमोप्रोपेन और मैग्नीशियम से शुष्क ईथर में \ce{(CH3)2CHMgBr} बनाइए,
उसे ठोस \ce{CO2} पर उँडेलिए, फिर अम्लीकृत कीजिए: \ce{(CH3)2CHCOOH}। $n = 12.3/123.0 = 0.100$ mol;
$0.100 \times 0.70 \times 88.0 = \qty{6.2}{g}$।
\end{solution}

\begin{solution}{exo:b1:organomagnesium:8}
$0.18/18.0 = \qty{0.010}{mol}$ जल \qty{0.010}{mol}
अभिकर्मक नष्ट करता है: \qty{20}{\percent}; बेंज़ीन बनती है, \ce{C6H5MgBr + H2O -> C6H6 +
Mg(OH)Br}।
\end{solution}

\begin{solution}{exo:b1:organomagnesium:9}
अभिकर्मक, जो एक \omterm{def:b1:electronic-effects:nucleophile}{नाभिकरागी} है, अब भी उपस्थित ब्रोमोबेंज़ीन से अभिक्रिया करता है:
कुल मिलाकर \ce{C6H5MgBr + C6H5Br -> C6H5-C6H5 + MgBr2}। ब्रोमाइड को
आधिक्य मैग्नीशियम में धीरे-धीरे मिलाने से उसकी सांद्रता कम रहती है: वह पहले से बने
अभिकर्मक के बजाय ताज़ी धातु से मिलता है।
\end{solution}

\begin{solution}{exo:b1:organomagnesium:10}
1-फ़ेनिलएथेन-1-ऑल: बेंज़ैल्डिहाइड के साथ \ce{CH3MgBr} (ब्रोमोमेथेन से),
या एथेनल के साथ \ce{C6H5MgBr} (ब्रोमोबेंज़ीन से)। 2-फ़ेनिलप्रोपेन-2-ऑल:
प्रोपेनोन के साथ \ce{C6H5MgBr}, या
1-फ़ेनिलएथेन-1-ओन के साथ \ce{CH3MgBr}।
\end{solution}

\begin{solution}{exo:b1:organomagnesium:11}
कीटोन को अभिकर्मक में मिलाया जाता है ताकि अभिकर्मक, जो शुष्क
फ़्लास्क में बना है, उसे कभी न छोड़े और हर भाग से मुक्त ऊष्मा
उबलते ईथर द्वारा अवशोषित हो जाए। केवल जल जिलेटिनी
क्षारकीय लवण \ce{Mg(OH)Br} देता, जो उत्पाद को फँसा लेता है; ठंडा तनु अम्ल
उसे \ce{Mg^2+} के रूप में घोल देता है, और ठंडक अम्ल में तृतीयक ऐल्कोहॉल की
पार्श्व अभिक्रियाओं को सीमित करती है।
\end{solution}

\begin{solution}{exo:b1:organomagnesium:12}
अभिकर्मक का हर मोल जल के साथ क्षारकीय मैग्नीशियम लवण का एक मोल देता है,
\ce{RMgBr + H2O -> RH + Mg(OH)Br}, जिसके \ce{OH-} का अम्ल से \omterm{def:b1:titration-methods:titration}{अनुमापन} होता है:
अभिकर्मक की मात्रा प्रयुक्त अम्ल के बराबर है, \qty{0.088}{mol}, अर्थात्
\qty{88}{\percent} \omterm{def:b1:extent-q-and-k:fractional-extent}{लब्धि}।
\end{solution}

\begin{solution}{pb:b1:organomagnesium:1}
\textbf{1.} \ce{C6H5Br + Mg -> C6H5MgBr}।
\textbf{2.} Mg $0.535/24.3 = \qty{0.0220}{mol}$; \ce{C6H5Br} $3.14/156.9 =
\qty{0.0200}{mol}$: मैग्नीशियम थोड़े आधिक्य में।
\textbf{3.} जल अभिकर्मक को नष्ट करता है: \ce{C6H5MgBr + H2O -> C6H6 +
Mg(OH)Br}।
\textbf{4.} यह वायु की जलवाष्प को संघनित्र से होकर भीतर आने से
रोकती है।
\textbf{5.} ईथर मैग्नीशियम को \omterm{def:b1:lewis-resonance-vsepr:lewis}{एकाकी युग्म} देता है (\omterm{def:b1:lewis-resonance-vsepr:lewis-acid}{लुइस क्षारक}):
उसके बिना अभिकर्मक न बनता है, न विलयन में रहता है।
\textbf{6.} \ce{C6H5MgBr + C6H5Br -> C6H5-C6H5 + MgBr2}; बाइफ़ेनिल
$0.15/154.0 = \qty{9.7e-4}{mol}$, जिसने \qty{9.7e-4}{mol}
ब्रोमोबेंज़ीन और उतना ही अभिकर्मक खपाया: कुल \qty{1.9e-3}{mol} ब्रोमोबेंज़ीन
की हानि।
\textbf{7.} $3.28/182.0 = \qty{0.0180}{mol}$।
\textbf{8.} C--Mg युग्म कार्बोनिल कार्बन पर आक्रमण करता है, $\pi$ युग्म
ऑक्सीजन पर जाता है: \ce{(C6H5)3C-O^-} \ce{MgBr+}।
\textbf{9.} उपलब्ध अभिकर्मक अधिक से अधिक $0.0200 - 0.0019 = \qty{0.0181}{mol}$;
बेंज़ोफ़ीनोन \qty{0.0180}{mol}: बेंज़ोफ़ीनोन सीमित करता है, बस थोड़े से।
\textbf{10.} नहीं: कार्बिनॉल कार्बन पर तीन सर्वसम फ़ेनिल समूह हैं।
\textbf{11.} अभिकर्मक अपने शुष्क फ़्लास्क में रहता है, और कीटोन के हर भाग की
ऊष्मा उबलता ईथर ले लेता है।
\textbf{12.} वह अभिकर्मक का कुछ भाग नष्ट कर देगा (बेंज़ीन) और
बेंज़ोफ़ीनोन को अनभिकृत छोड़ देगा।
\textbf{13.} \ce{(C6H5)3COMgBr + H3O+ -> (C6H5)3COH + Mg^2+ + Br- + H2O}।
\textbf{14.} अम्ल उन क्षारकीय मैग्नीशियम लवणों को घोल देता है जिन्हें केवल जल
अवक्षेपित कर देता; ठंडक प्रश्न 18 के
\omterm{def:b1:electronic-effects:intermediates}{कार्बधनायन} के बनने को सीमित करती है।
\textbf{15.} ट्राइफ़ेनिलमेथेनॉल और बाइफ़ेनिल ईथर परत में;
मैग्नीशियम लवण जलीय परत में।
\textbf{16.} ठोस को ठंडे हल्के पेट्रोलियम से धोकर (या रगड़कर):
बाइफ़ेनिल घुल जाता है, ट्राइफ़ेनिलमेथेनॉल बचा रहता है।
\textbf{17.} सारणीबद्ध मान के बराबर एक तीक्ष्ण गलनांक (या पतली परत वर्णलेखन में
एक ही धब्बा)।
\textbf{18.} \ce{(C6H5)3COH + H+ -> (C6H5)3C+ + H2O}: एक तृतीयक \omterm{def:b1:electronic-effects:intermediates}{कार्बधनायन}
जिसका आवेश अनुनाद द्वारा तीन बेंज़ीन वलयों पर फैला है, अत्यंत
स्थायी, और उसी विस्तृत संयुग्मन के कारण रंगीन।
\textbf{19.} $3.75/260.0 = \qty{0.0144}{mol}$।
\textbf{20.} $0.0144/0.0180 = \qty{80}{\percent}$।
\textbf{21.} स्थानांतरणों और पुनःक्रिस्टलन में हानियाँ; नमी के अंशों से नष्ट हुआ
अभिकर्मक; अपूर्ण अभिक्रिया।
\textbf{22.} अम्लीकरण के बाद बेंज़ोइक अम्ल, \ce{C6H5COOH}।
\textbf{23.} $0.0181 \times 122.0 = \qty{2.21}{g}$।
\textbf{24.} \ce{C6H5MgBr + CO2 -> C6H5COOMgBr}, फिर
\ce{C6H5COOMgBr + H3O+ -> C6H5COOH + Mg^2+ + Br- + H2O}।
\textbf{25.} \textbf{\qty{80}{\percent}}।
\end{solution}
